% !TeX root = ../se320_structural_assurability.tex

% ==================================================
\section{Problem Setting}
\label{aa.sec.foundations}
% ==================================================

Assurance is fundamentally concerned with claims and evidence.
Existing assurance theory provides substantial machinery for reasoning
from evidence to claims: whether evidence is relevant and sufficient,
whether an assurance argument is sound, what assumptions it depends upon,
what defeaters remain, and what confidence in the resulting claim is justified.

This work considers an antecedent question:

\begin{quote}
  What structural properties determine whether evidence relevant to
  supporting or refuting an assurance claim can be obtained in the first place?
\end{quote}

This question is distinct from whether the claim is true and from whether
a particular assurance argument is convincing.
A system may satisfy a desired property while exposing little evidence
that it does so.
Conversely, a system may expose extensive and reliable evidence
demonstrating that a desired property is not satisfied.

We use the term \emph{assurability} for the capacity to make
assurance-relevant evidence obtainable.
We use \emph{structural assurability} when emphasizing architectural
and environmental properties that contribute to this capacity.

The objective of this section is deliberately limited.
We define the object of study, identify the context in which assurability
must be evaluated, distinguish structural properties from evidence-quality
properties and practical constraints, and identify candidate structural
dimensions whose contribution can be demonstrated through separating
constructions.

We do not claim that these dimensions are jointly sufficient to
characterize assurability, individually necessary for every assurance
problem, mutually independent, or complete.

\subsection{From Assurance Evidence to Evidence Availability}
\label{aa.subsec.evidence-availability}

Let \(S\) denote a system and \(C\) an assurance claim concerning \(S\).
Let \(E\) denote evidence relevant to \(C\), and let \(G\) denote an
assurance argument connecting that evidence, together with applicable
assumptions, to the claim.

At a high level,

\begin{equation}
  E \xrightarrow{G} C.
  \label{eq:evidence-claim}
\end{equation}

Much assurance research appropriately concentrates on this relationship:
whether \(E\) supports \(C\), whether \(G\) is sound, what counterevidence
or defeaters exist, and what confidence is justified.

Before such reasoning can occur, however, relevant evidence must be obtainable.

Consider systems \(S_1\) and \(S_2\) that implement the same externally
visible functionality and exhibit equivalent behavior over an evaluation
set \(T\):

\begin{equation}
  \operatorname{Behavior}_{T}(S_1)
  =
  \operatorname{Behavior}_{T}(S_2).
  \label{eq:behavior-equivalence}
\end{equation}

Suppose \(S_1\) exposes only externally visible inputs and outputs,
while \(S_2\) additionally preserves authenticated records of
consequential actions, tool invocations, retrieved information,
authorization decisions, provenance, and execution state.

Behavioral equivalence does not imply evidentiary equivalence.
The systems may therefore be indistinguishable with respect to \(T\)
while differing materially in what can subsequently be established
about their operation.

This observation has mature analogues in other engineering disciplines.
Design for testability modifies system architecture to improve the ability
to establish test conditions and observe relevant system behavior.
Software architecture similarly affects testability, analyzability,
monitorability, and the cost and effectiveness of verification activities
\cite{abramovici1990digitaltesting,bass2021softwarearchitecture}.

The corresponding structural principle for assurance is:

\begin{quote}
  The architecture of a system and its evidence-producing environment
  constrain the evidence that can practically be obtained concerning
  assurance claims about that system.
\end{quote}

\subsection{Definition of Assurability}
\label{aa.subsec.definition-assurability}

The preceding observation motivates a deliberately minimal definition.

\begin{definition}[Assurability]
  \label{se320.def.Assurability}
  The \emph{assurability} of a system \(S\) with respect to an assurance
  claim \(C\), under conditions \(K\) and resource constraints \(B\),
  is the capacity of the system and its evidence-producing environment
  to make evidence relevant to supporting or refuting \(C\)
  practically obtainable.
\end{definition}

We denote this claim-relative capacity by

\begin{equation}
  \mathcal{A}(S,C,K,B).
  \label{eq:assurability}
\end{equation}

The notation does not imply that \(\mathcal{A}\) is scalar-valued.
Indeed, the framework developed here deliberately avoids assuming
that heterogeneous assurance properties can be reduced to a meaningful
universal score.

Several features of Definition~1 are intentional.

First, assurability concerns evidence, not truth.
A highly assurable system may provide decisive evidence that claim \(C\)
is false.

Second, assurability concerns evidence capable of supporting
\emph{or refuting} \(C\).
A structure that facilitates favorable evidence while systematically
obscuring disconfirming evidence should not thereby be considered
highly assurable.

Third, the definition does not specify which evidence is sufficient
for establishing \(C\).
Sufficiency belongs to the downstream assurance problem and depends
upon the claim, argument, evidence model, threat model, and applicable
assurance requirements.

Fourth, the definition does not require that evidence originate within
the system.
The evidence-producing environment may include deployment infrastructure,
independent monitors, external evaluators, physical sensors, network
services, human processes, or other systems.

Assurability is therefore distinct from the property being assured,
from the assurance argument, and from confidence in that argument:

\begin{equation}
  \boxed{
    \begin{aligned}
      \text{Property satisfaction}
       & \neq
      \text{Assurability} \\
       & \neq
      \text{Assurance argument}
      \neq
      \text{Confidence in the argument}.
    \end{aligned}
  }
  \label{eq:assurance-distinctions}
\end{equation}

\subsection{Assurance Context}
\label{aa.subsec.assurance-context}

Assurability is not an unconditional property of \(S\).

A system may provide strong evidence concerning whether an external
action occurred while providing little evidence concerning the internal
computational process that produced the action.
It may therefore be highly assurable with respect to one claim and
weakly assurable with respect to another.

Evidence availability also depends upon evaluator access and cooperation.
A developer may possess model weights, activations, system prompts,
execution traces, configuration data, and infrastructure logs.
An external evaluator may possess only API access.
An independent assessor may receive selected evidence but lack direct
system access.
In an adversarial setting, the developer, operator, or system itself
may attempt to suppress, manipulate, or fabricate evidence.

We collect these knowledge, access, trust, and cooperation conditions
in \(K\).

Finally, theoretically obtainable evidence may not be practically
obtainable.
Assurance consumes computation, time, expertise, storage, instrumentation,
system access, and other resources.
Assurance mechanisms may additionally impose latency, performance,
disclosure, or operational costs.
We represent relevant resource constraints by \(B\).

We therefore define:

\begin{definition}[Assurance Context]
  \label{se320.def.AssuranceContext}
  An \emph{assurance context} is a tuple
  \begin{equation}
    Q=(S,C,K,B),
  \end{equation}
  specifying the system under consideration, the assurance claim,
  the applicable knowledge, access, trust, and cooperation conditions,
  and the resource constraints under which assurance evidence
  must be obtained.
\end{definition}

Comparative statements about assurability are meaningful only relative
to a sufficiently specified assurance context.

In particular,

\begin{quote}
  System \(S_1\) is more assurable than system \(S_2\)
\end{quote}

is generally underspecified.

The comparison must identify what is being assured and under what conditions.

\subsection{Structure, Evidence Quality, and Constraints}
\label{aa.subsec.three-layers}

A useful characterization of assurability requires distinguishing
three classes of factors that are easily conflated:

\begin{enumerate}
  \item structural properties of the system and evidence-producing environment;
  \item properties governing the evidentiary or inferential quality of observations; and
  \item contextual constraints governing whether evidence can practically be obtained.
\end{enumerate}

These occupy different roles in the assurance process.

\paragraph{Structural properties.}
Structural properties determine what can be exposed, controlled,
associated with its origin, preserved, reconstructed, independently
corroborated, protected, and covered by the evidence-producing architecture.
These properties are the primary subject of structural assurability.

\paragraph{Evidence and inference properties.}
Once evidence is available, additional questions arise concerning its
relevance, validity, fidelity, calibration, and uncertainty.
These properties govern what may legitimately be inferred from an
observation or measurement.

For example, an architecture may expose an internal measurement perfectly
while the measurement itself is an invalid indicator of the construct
it is intended to represent.
Greater structural observability does not resolve construct validity.

\paragraph{Contextual constraints.}
Access, cooperation, time, computation, expertise, financial cost,
disclosure restrictions, and operational burden constrain which
structurally available evidence can actually be obtained in a
particular assurance context.
These constraints are represented principally through \(K\) and \(B\).

The distinction can be summarized as

\begin{equation}
  \boxed{
    \begin{aligned}
      \text{Structure}
       & \longrightarrow
      \text{Obtainable Evidence} \\
       & \longrightarrow
      \text{Evidence and Inference}
      \longrightarrow
      \text{Assurance Argument},
    \end{aligned}
  }
  \label{eq:three-layer-chain}
\end{equation}

subject throughout to the conditions and constraints of the assurance context.

This separation is important because the present theory concerns
the first transition.
Structural assurability constrains what evidence can be made available;
it does not by itself establish the validity of subsequent inference
or the sufficiency of an assurance argument.

\subsection{Candidate Structural Dimensions}
\label{aa.subsec.candidate-dimensions}

Prior work across hardware testing, software architecture, safety
engineering, security evaluation, auditability, runtime assurance,
and AI monitoring suggests several recurring structural properties
that affect the availability and usefulness of assurance evidence.

We begin with seven candidate dimensions:

\begin{description}
  \item[Observability.]
        The extent to which states, events, behaviors, or effects relevant
        to an assurance claim can be observed.

  \item[Controllability.]
        The extent to which conditions relevant to evaluating an assurance
        claim can be established and relevant interventions can be performed.

  \item[Traceability.]
        The extent to which evidence can be associated with its origin,
        transformation history, and relevant system events.

  \item[Reconstructability.]
        The extent to which a relevant execution, decision sequence,
        or system state can subsequently be reconstructed from preserved evidence.

  \item[Independence.]
        The extent to which evidence can be generated, preserved, or
        corroborated independently of the component or actor whose behavior
        is the subject of the assurance claim.

  \item[Integrity.]
        The extent to which alteration, fabrication, suppression,
        substitution, or corruption of assurance evidence can be detected
        or prevented.

  \item[Evidence Coverage.]
        The extent to which the evidence-producing structure exposes
        the states, behaviors, events, interfaces, or effects material
        to the assurance claim.
\end{description}

We denote the current candidate set by

\begin{equation}
  D_S =
  \{
  O,C,T,R,N,I,E
  \},
  \label{eq:structural-dimensions}
\end{equation}

where the symbols are used only as compact identifiers and do not
imply numerical measurement scales.

The candidate set is deliberately provisional.

The dimensions are not asserted to be mutually independent.
They may interact or partially depend upon one another.
Reconstructability, for example, may depend upon observability,
traceability, and evidence integrity.
Evidence coverage may depend upon the placement and scope
of observation mechanisms.

Nor is each dimension asserted to be necessary for every assurance claim.
Some claims require no access to internal state.
Some require no intervention.
Some can be established through independently observable external effects.

The purpose of the candidate dimensions is therefore not to define
a checklist that every assurable system must satisfy.
It is to identify structural variables whose omission can erase
distinctions that matter to assurance.

\subsection{Derived Assurance Capabilities}
\label{aa.subsec.derived-capabilities}

Several important properties associated with assurance appear better
treated as capabilities emerging from combinations of structural
dimensions rather than as primitive dimensions themselves.

These include, potentially,

\begin{equation}
  \{
  \text{reproducibility},
  \text{auditability},
  \text{monitorability},
  \text{testability},
  \text{diagnosability}
  \}.
\end{equation}

For example, reproducibility may depend upon sufficient traceability,
reconstructability, controllability, configuration preservation,
and environmental information.
Testability classically depends in substantial part upon observability
and controllability.

Treating these properties as derived rather than primitive avoids
assuming independence where the underlying engineering relationships
suggest composition.

This classification is itself provisional.
If a proposed derived capability contains an assurance-relevant
structural property not represented by the current primitive set,
that property should be identified and tested separately.

\subsection{The Dimension Contribution Criterion}
\label{aa.subsec.dimension-contribution}

A candidate structural dimension should not be included merely
because it appears intuitively related to assurance.
Its contribution should be demonstrable.

We therefore introduce a criterion based on separating constructions.

\begin{definition}[Assurance-Relevant Structural Dimension]
  \label{se320.def.AssuranceRelevantStructuralDimension}
  A candidate structural dimension \(d\) is \emph{assurance-relevant}
  if there exists an assurance context and a pair of systems for which
  a difference in \(d\), while holding other properties material to
  the comparison fixed, produces a material difference in the evidence
  practically obtainable for supporting or refuting the assurance claim.
\end{definition}

Let \(S_d^{+}\) and \(S_d^{-}\) denote systems constructed to differ
materially with respect to candidate dimension \(d\) while remaining
equivalent with respect to other properties material to a specified comparison.

Let \(\mathcal{E}(S,C,K,B)\) denote the claim-relevant evidence
practically obtainable from \(S\) under conditions \(K\)
and constraints \(B\).

Dimension \(d\) is assurance-relevant if there exist \(C\), \(K\),
and \(B\) such that

\begin{equation}
  \mathcal{E}(S_d^{+},C,K,B)
  \not\equiv_C
  \mathcal{E}(S_d^{-},C,K,B),
  \label{eq:dimension-separation}
\end{equation}

where \(\not\equiv_C\) denotes a difference in obtainable evidence
material to supporting or refuting \(C\).

We call such a pair a \emph{separating construction} for \(d\).

The logical strength of this criterion is intentionally limited.

A separating construction does not establish that \(d\) is required
for every assurance context.
It does not establish that increasing \(d\) monotonically increases
assurability.
It does not establish that \(d\) is independent of other dimensions.

It establishes only that

\begin{quote}
  there exists an assurance context in which omitting \(d\) from
  the model would collapse a distinction between systems that
  materially affects their capacity to support or refute an
  assurance claim.
\end{quote}

This is the sense in which a dimension is necessary to the general
framework: not because every assurable system must possess it,
but because a framework incapable of representing it would fail
to distinguish some systems that differ materially in assurability.

\subsection{Separating the Candidate Dimensions}
\label{aa.subsec.separating-dimensions}

The contribution criterion can be illustrated informally for each
current candidate dimension.

\paragraph{Observability.}
Consider otherwise equivalent systems in which only one exposes
a state relevant to \(C\).
If direct observation of that state provides evidence unavailable
from the other system, observability separates their assurability.


\paragraph{Controllability.}
Consider otherwise equivalent observable systems in which only one
permits the evaluator to establish a condition or intervention
required to investigate \(C\).
The ability to obtain evidence under that condition separates
their assurability.

\paragraph{Traceability.}
Consider systems that expose the same event but differ in whether
the event can be reliably associated with its originating component,
action, or causal history.
For a claim whose evaluation depends upon attribution, traceability
separates their assurability.

\paragraph{Reconstructability.}
Consider systems with equivalent runtime behavior and live
observability, but only one preserves sufficient state and history
to reconstruct a relevant execution after an incident.
For retrospective assurance claims, reconstructability separates
their assurability.

\paragraph{Independence.}
Consider systems providing equivalent records, but in one system
the only record is generated by the component being evaluated,
while the other additionally provides independently generated
or protected corroboration.
For claims vulnerable to self-reporting or common-mode failure,
independence separates their assurability.

\paragraph{Integrity.}
Consider systems exposing equivalent records, but only one
prevents or detects alteration, fabrication, or suppression
of those records.
For claims relying upon the historical record, integrity
separates their assurability.

\paragraph{Evidence Coverage.}
Consider systems whose individual observations are otherwise
equivalent, but one exposes only a proper subset of the
claim-relevant action surface.
For claims spanning the omitted behavior, evidence coverage
separates their assurability.
These examples motivate the dimension-specific constructions
examined in Section~\ref{aa.sec.formal-model}.

Section~\ref{aa.sec.formal-model} develops the comparative relation
required to state these results more precisely.

\subsection{Contribution Without Sufficiency}
\label{aa.subsec.necessity-sufficiency}

The framework deliberately distinguishes between demonstrating
that a dimension contributes to assurability and claiming that
a collection of dimensions completely characterizes assurability.

For a candidate dimension \(d\), a separating construction can
demonstrate that \(d\) cannot generally be discarded without
losing an assurance-relevant distinction.

It does not establish that \(d\) is necessary for every assurance claim.

Likewise, establishing the relevance of every member of the current
candidate set \(D_S\) does not establish that \(D_S\) is sufficient
or complete.

In particular, we do not claim

\begin{equation}
  \bigwedge_{d_i \in D_S} d_i
  \Longrightarrow
  \text{sufficient assurance},
  \label{eq:no-sufficiency}
\end{equation}

nor do we claim that \(D_S\) exhausts the structural properties
relevant to future assurance problems.

A new system architecture, assurance claim, threat model,
measurement mechanism, or evidence source may expose an
additional structural property \(d_{n+1}\).

If that property admits a separating construction under
Definition~3, the framework can be extended:

\begin{equation}
  D'_S = D_S \cup \{d_{n+1}\}.
\end{equation}

This open-endedness is intentional.

The space of future assurance claims, architectures, threats,
measurements, and evidence mechanisms is not closed.
A claim that the current dimensions are jointly sufficient
would therefore require assumptions that this work cannot justify.

The contribution is narrower:

\begin{quote}
  We identify structural dimensions that demonstrably affect
  assurability, provide a criterion for establishing their
  contribution, and construct a framework that can be extended
  when additional assurance-relevant dimensions are identified.
\end{quote}

\subsection{Structural Neutrality}
\label{aa.subsec.structural-neutrality}

Structural assurability does not determine the outcome of assurance.

An evidence-producing architecture should support examination
regardless of whether the resulting evidence confirms or defeats
the assurance claim.
The structure should not embed the favorable conclusion
it is intended to help evaluate.

Thus,

\begin{equation}
  \text{greater structural capacity for assurance}
  \not\Rightarrow
  \text{a more favorable assurance outcome}.
\end{equation}

Indeed, a highly assurable system may be easier to reject precisely
because its structure makes adverse evidence difficult to conceal.

This neutrality separates evidence availability from downstream
interpretation and judgment.
Structural assurability concerns the conditions under which relevant
evidence can be obtained and examined; it does not determine what
conclusion an evaluator must draw from that evidence.

\subsection{The Comparative Problem}
\label{aa.subsec.comparative-problem}

The preceding definitions lead to the central problem of comparative
assurability.

Consider systems \(S_1\) and \(S_2\) subject to the same assurance
claim \(C\), knowledge and access conditions \(K\), and resource
constraints \(B\).

The systems may exhibit equivalent functional behavior while differing
along one or more assurance-relevant structural dimensions.

The comparative question is:

\begin{quote}
  Do structural differences between \(S_1\) and \(S_2\) create
  material differences in the evidence practically obtainable
  for supporting or refuting \(C\)?
\end{quote}

\begin{definition}[Evidentiary Equivalence]
  \label{se320.def.EvidentiaryEquivalence}
  Two systems \(S_1\) and \(S_2\) are \emph{evidentially equivalent}
  relative to a claim \(C\), conditions \(K\), and constraints \(B\)
  if their practically obtainable evidence has no difference
  material to supporting or refuting \(C\). We write
  \[
    \mathcal{E}(S_1,C,K,B)
    \equiv_C
    \mathcal{E}(S_2,C,K,B).
  \]
\end{definition}

\begin{definition}[Comparative Assurability]
  \label{se320.def.ComparativeAssurability}
  \emph{Comparative assurability} evaluates the relative capacity
  of two systems to make evidence relevant to supporting or
  refuting the same claim practically obtainable under specified,
  comparable conditions and resource constraints.
  The comparison may yield equivalence, dominance, tradeoffs,
  or incomparability; it does not presuppose a scalar score
  or total order.
\end{definition}

Behavioral equivalence does not imply assurability equivalence.

This motivates the central premise of the framework:

\begin{quote}
  \textbf{Central Premise.}

  Systems with equivalent functional behavior or measured
  performance can differ materially in assurability because
  their architectures and evidence-producing environments
  can differ in assurance-relevant structural dimensions.
\end{quote}

The contribution is therefore neither a universal assurance score
nor a finite checklist whose satisfaction guarantees assurance.

It is a framework for identifying, demonstrating, and comparing
structural factors that affect the practical capacity to obtain
assurance evidence.

Section~\ref{aa.sec.formal-model} develops the comparative relation
formally and examines the candidate dimensions through separating
constructions, while permitting equivalence, dominance, tradeoffs,
and incomparability rather than requiring a total ordering.
